\documentclass[12pt]{article}
\usepackage[letterpaper,margin=0.75in]{geometry}
\usepackage[T1]{fontenc}
\usepackage{lmodern}
\usepackage{microtype}
\usepackage{booktabs}
\usepackage{array}
\usepackage{hyperref}
\usepackage{setspace}
\setstretch{1.0}
\setlength{\parindent}{1em}
\setlength{\parskip}{0pt}
\setlength{\tabcolsep}{5pt}
\hypersetup{hidelinks,pdftitle={Romanization or resolution? A shuffled-label control for Cherokee-English translation memory},pdfauthor={},pdfsubject={Anonymous PLC 51 abstract},pdfkeywords={}}
\pagestyle{plain}
\renewcommand{\refname}{References}

\begin{document}

\begin{center}
{\bfseries Romanization or resolution? A shuffled-label control for\\
Cherokee--English translation memory}
\end{center}
\vspace{-0.65em}

Transliteration can improve low-resource language technology by exposing
cross-script or phonological overlap (Goyal et al. 2020; Moosa et al. 2023).
But a raw-script--romanized comparison changes two things at once. It changes
the linguistic labels, and it expands one source symbol into several Latin
characters. The latter alone alters the geometry of character n-gram models.
We isolate these factors in Cherokee--English translation, where the Cherokee
syllabary provides an unusually clean test: a source symbol has a conventional
Latin syllable label, but it can also be assigned an arbitrary label of the
same length. A public corpus experiment finds a significant $+1.02$ chrF2
gain from romanization over raw syllabary. Crucially, length-matched shuffled
labels reproduce about 80\% of that gain. Thus the usual two-condition
comparison overstates the evidence that transliteration helps because it
encodes linguistically appropriate sub-symbol structure.

\textbf{Background.}
Cherokee is polysynthetic and severely endangered. ChrEn provides 14K
Cherokee--English pairs and separate in- and out-of-domain evaluations
(Zhang et al. 2020). Work with Cherokee speakers also emphasizes that language
technology must be developed through community partnership rather than treating
an endangered language as a benchmark alone (Zhang et al. 2022). Our target is
therefore methodological: we audit a common modeling inference, not propose a
new community orthography. The romanization below is an automatically
reproducible diagnostic based on Unicode character labels.

\textbf{Method.}
We use the published ChrEn splits: 11,639 training, 1,000 development, 1,000
in-domain test, and 256 out-of-domain test pairs. To make the representational
effect transparent, the translation system is one-nearest-neighbor translation
memory. A TF--IDF vector represents each Cherokee sentence with character
$n$-grams; the English side of its nearest training neighbor is returned as the
translation. Native-script and romanized $n$-gram ranges are independently
selected on development data from 1--4, 1--6, 2--5, 2--6, and 2--8. Both select
1--4. Output is scored with corpus chrF2 over character 1--6-grams
(Popovic 2015).

The \textsc{raw} condition uses Cherokee syllabary characters. The
\textsc{romanized} condition maps each of 85 Cherokee symbols losslessly to the
lowercase label in its Unicode character name (for example, a syllabary symbol
labeled \emph{NA} becomes \texttt{na}). This exposes within-syllable Latin
substrings to the retriever. The key control randomly permutes these 85 labels
among symbols \emph{within label-length groups}. Each control therefore
preserves the romanized alphabet, label frequencies, and every symbol's
expansion length, while destroying the association between a symbol and its
intended syllable label. We run 20 permutations. Uncertainty for the raw versus
romanized difference comes from 2,000 paired bootstrap samples of test items.

\textbf{Results.}
Table~\ref{tab:results} shows that romanization raises in-domain chrF2 from
34.74 to 35.76, a $+1.02$ gain whose bootstrap 95\% interval is
$[0.49,1.51]$; none of 2,000 resamples is non-positive. Source-side nearest-
neighbor cosine similarity also rises from .344 to .443. The apparent benefit
therefore replicates what an ordinary raw--romanized comparison would call a
successful transliteration intervention.

\begin{table}[h]
\centering
\caption{Cherokee-to-English translation-memory retrieval.}
\label{tab:results}
\vspace{-0.1em}
\begin{tabular}{@{}lrrr@{}}
\toprule
Representation & In-domain chrF2 & Out-domain chrF2 & Source cosine\\
\midrule
Raw syllabary & 34.74 & 18.54 & .344\\
Unicode-label romanization & 35.76 & 19.21 & .443\\
Shuffled labels (20; mean $\pm$ SD) & $35.55\pm.15$ & -- & --\\
\bottomrule
\end{tabular}
\end{table}
\vspace{-0.6em}

The shuffled-label result changes the interpretation. Its scores range from
35.39 to 35.92; two of 20 arbitrary mappings outperform the linguistically
labeled mapping. The romanized condition is only 0.21 above the shuffle mean,
while the shuffle mean is 0.81 above raw. On this operationalization, roughly
80\% of the headline gain is recoverable without correct syllable labels.

\textbf{Analysis.}
A Cherokee syllabary character is atomic to the raw character model. Expanding
it into multiple Latin characters creates additional partial matches, longer
contexts, and more gradations of similarity. A random length-matched label does
the same. Consequently, a raw--romanized contrast conflates a linguistic claim
(the correct sub-symbol correspondence matters) with a resolution claim (more
character positions help). The shuffle test supplies the missing counterfactual.
The 0.21 advantage over the control mean, and the fact that 90\% of controls
score lower, leave room for a small linguistically structured effect, but 20
permutations do not distinguish it decisively.

Romanization also does not remove the corpus's domain problem. Both systems
lose about 16 chrF2 points out of domain; the romanized advantage there is only
0.67. This echoes the original ChrEn finding that out-of-domain Cherokee
translation is substantially harder. More broadly, transliteration is not a
theory-neutral preprocessing step. It chooses which internal similarities a
model may see and can privilege segmental Latin structure over syllabary-based
orthographic practice. Claims of phonological or cross-script benefit should
therefore include matched nonlinguistic recodings, report domain transfer, and
ultimately be validated with speaker-approved representations and human
judgment. The proposed control is cheap, model-agnostic, and applicable to any
low-resource setting in which transliteration expands source symbols.

\vspace{-0.3em}
\begin{thebibliography}{9}
\setlength{\itemsep}{0pt}
\setlength{\parskip}{0pt}

\bibitem[Goyal et al. 2020]{goyal2020}
Goyal, Vikrant, Sourav Kumar, and Dipti Misra Sharma. 2020.
\newblock Efficient neural machine translation for low-resource languages via exploiting related languages.
\newblock \textit{ACL Student Research Workshop}. 162--168.

\bibitem[Moosa et al. 2023]{moosa2023}
Moosa, Ibraheem Muhammad, Mahmud Elahi Akhter, and Ashfia Binte Habib. 2023.
\newblock Does transliteration help multilingual language modeling?
\newblock \textit{Findings of EACL}. 670--685.

\bibitem[Popovic 2015]{popovic2015}
Popovic, Maja. 2015.
\newblock chrF: Character n-gram F-score for automatic MT evaluation.
\newblock \textit{WMT}. 392--395.

\bibitem[Zhang et al. 2020]{zhang2020}
Zhang, Shiyue, Benjamin Frey, and Mohit Bansal. 2020.
\newblock ChrEn: Cherokee--English machine translation for endangered language revitalization.
\newblock \textit{EMNLP}. 577--595.

\bibitem[Zhang et al. 2022]{zhang2022}
Zhang, Shiyue, Ben Frey, and Mohit Bansal. 2022.
\newblock How can NLP help revitalize endangered languages? A case study and roadmap for Cherokee.
\newblock \textit{ACL}. 1529--1541.

\end{thebibliography}

\end{document}
